\scanpage{90}
% loc 1–4, upper display and gloss
\[
\left\{\begin{aligned}
\rho_0=\rho-2
&=\frac{\delta\deleted{(u)}\delta'\deleted{(v)}}{t^2-bb't+(\delta c'+\delta'c+4\deleted{\unsure{b}}cc')}\\
&=\frac{(b^2-4c)(b'^2-4c')}{\underbrace{t^2-bb't+b^2c'+b'^2c-4cc'}_{Q=Q(t;b,b',c,c')}}
\end{aligned}\right.\qquad
\begin{aligned}\text{où }\delta&=\delta(u)=b^2-4c\\ \delta'&=\delta(v)=b'^2-4c'\end{aligned}
\tag*{$(11')$}
\]
% loc 5
\unsure{$1$}

On trouve donc :

% loc 6–7, proposition
\uline{Proposition} Soient $k$ un anneau, et $b,c,b',c',t\in k$, tels que
% loc 8
\[ (b^2-4c)(b'^2-4c')(t^2-bb't+b^2c'+b'^2c-4cc') \]
% loc 9–15
soit inv. dans $k$. Alors $M[b,c,b',c',t]$ est une alg. d'Azumaya. (Et $k[\underline u],k[\underline v]$ en sont des sous-algèbres quadratiques étales, \inserted{``en position générale'' i.e \unsure{engendrent|engendrant} \unsure{la dite|de cette} Alg.}) Dans tous les cas, c'est une algèbre associative, \deleted{et} satisfaisant les relations (2) (3) pour $\operatorname{Tr},\det$ définis comme plus haut\ldots.

% loc 16–20, lower-left marginal OBS
% slanted
\margin{OPS $k$ corps alg. clos, car c'est une question \unsure{géom.} fibre par fibre.}
% loc 16–21, proof
\uline{Dém} On regarde $\rho_0$ comme défini plus haut $(11')$. Comme $\rho_0$ \deleted{\unsure{étant}} inv i.e $\rho$ partout $\ne2$, cela correspond à une position relative de tores dans un $\operatorname{SL}(2)$ (ou de sous alg. quadratiques \inserted{$A_1,A_2$} étales dans un $M_2$) en \uline{position générale}. -- Cela signifie aussi que
% archive sheet number: 89
